\chapter{PENDAHULUAN}
\label{chap:pendahuluan}

\section{Latar Belakang}
DBLP (\textit{Digital Bibliography \& Library Project}) adalah basis data bibliografi publikasi ilmiah bidang ilmu komputer. Setiap publikasi tercatat dengan judul, penulis, tahun, dan tempat terbitnya \cite{ley2009}. Riwayat publikasi para peneliti pada DBLP dapat dipakai untuk menjawab dua pertanyaan. Pertama, konferensi mana yang paling aktif. Kedua, konferensi mana yang sering dipilih oleh peneliti yang sama.

Pertanyaan kedua merupakan persoalan \textit{frequent pattern mining}. Jika setiap peneliti dipandang sebagai satu transaksi dan konferensi tempat ia memublikasikan karya dipandang sebagai item, kombinasi konferensi yang sering muncul bersama dapat dicari dengan algoritma Apriori \cite{agrawal1994} dan FP-Growth \cite{han2000}, lalu diturunkan menjadi \textit{association rules}. Pada laporan ini, popularitas konferensi diukur dari aktivitas publikasi, yaitu jumlah publikasi dan jumlah penulis unik.

\section{Tujuan}
Tujuan tugas besar ini adalah sebagai berikut.
\begin{enumerate}
\item Memahami struktur data DBLP dan melakukan \textit{preprocessing} terhadap publikasi konferensi periode 2015--2025 (\textit{Task} 1).
\item Menganalisis popularitas konferensi berdasarkan jumlah publikasi dan jumlah penulis unik (\textit{Task} 2).
\item Membentuk \textit{transactional database} berbasis penulis dan berbasis penulis-tahun (\textit{Task} 3).
\item Mengimplementasikan Apriori dan FP-Growth untuk menemukan \textit{frequent itemsets} (\textit{Task} 4).
\item Membangkitkan dan menganalisis \textit{association rules} antarkonferensi (\textit{Task} 5).
\item Membandingkan performa Apriori dan FP-Growth (\textit{Task} 6).
\item Menganalisis perubahan pola asosiasi antarperiode waktu (Bagian E).
\end{enumerate}

\section{Ruang Lingkup}
\begin{enumerate}
\item Data yang digunakan adalah \textit{XML dump} DBLP \cite{dblpxml} yang diunduh pada 19 November 2026.
\item Jenis publikasi yang dianalisis hanya \texttt{inproceedings}, yaitu makalah konferensi dan \textit{workshop}.
\item Periode publikasi yang dianalisis adalah 2015--2025.
\item Atribut yang digunakan adalah \textit{key}, \textit{author}, \textit{title}, \textit{year}, dan \textit{booktitle}.
\item Popularitas konferensi hanya diukur dari aktivitas publikasi, bukan dari sitasi atau kualitas.
\end{enumerate}

\section{Sistematika Penulisan}
\begin{enumerate}
\item Bab~\ref{chap:pendahuluan} Pendahuluan berisi latar belakang, tujuan, ruang lingkup, dan sistematika penulisan.
\item Bab~\ref{chap:dataset} Deskripsi \textit{Dataset} dan Dasar Teori berisi struktur data DBLP serta konsep \textit{frequent pattern mining} dan \textit{association rule mining}.
\item Bab~\ref{chap:metodologi} Metodologi berisi penerapan tahap CRISP-DM pada proyek ini.
\item Bab~\ref{chap:prepro} Hasil \textit{Preprocessing} berisi jenis publikasi, hasil tiap tahap, dan statistik \textit{dataset} bersih.
\item Bab~\ref{chap:popularitas} Hasil \textit{Conference Popularity Analysis} berisi peringkat, tren, perbandingan, dan uji sensitivitas.
\item Bab~\ref{chap:transaksi} Hasil Pembentukan \textit{Transactional Database} berisi dua representasi transaksi beserta perbandingannya.
\item Bab~\ref{chap:fpm} sampai Bab~\ref{chap:temporal} berisi hasil \textit{frequent pattern mining}, \textit{association rule mining}, perbandingan performa, dan analisis temporal.
\item Bab~\ref{chap:analisis} Analisis dan Interpretasi menafsirkan temuan yang tersedia.
\item Bab~\ref{chap:kesimpulan} Kesimpulan berisi jawaban pertanyaan kesimpulan yang sudah dapat dijawab, keterbatasan, dan asumsi.
\end{enumerate}
